\documentclass[11pt]{article}
\usepackage[T1]{fontenc}
\usepackage{lmodern,amsmath,amssymb,amsthm,mathtools,microtype}
\usepackage[margin=0.86in,headheight=15pt]{geometry}
\usepackage{booktabs,tabularx,array,xcolor,fancyhdr,xurl}
\usepackage[colorlinks=true,linkcolor=teal!65!black,urlcolor=teal!65!black]{hyperref}
\definecolor{ink}{HTML}{173246}
\definecolor{accent}{HTML}{006F73}
\pagestyle{fancy}\fancyhf{}\fancyhead[L]{\small\sffamily\color{ink}Direct module duality in Ribet's method}\fancyhead[R]{\small\sffamily Research proposal}\fancyfoot[R]{\small\sffamily\thepage}
\setlength{\parindent}{0pt}\setlength{\parskip}{5pt plus 1pt}\setlength{\emergencystretch}{2em}\raggedbottom
\newtheorem{theorem}{Theorem}[section]\newtheorem{lemma}[theorem]{Lemma}\newtheorem{proposition}[theorem]{Proposition}
\newcommand{\OO}{\mathcal O}\newcommand{\FF}{\operatorname{Fitt}_0}\newcommand{\len}{\operatorname{length}_{\OO}}\newcommand{\val}{\operatorname{v}}\newcommand{\Hom}{\operatorname{Hom}}\newcommand{\Span}{\operatorname{span}}\newcommand{\code}[1]{\texttt{\detokenize{#1}}}
\newcommand{\src}[2]{\href{https://github.com/VladimirReshetnikov/ProveIt/blob/9050b04000647d588d0c78d945180abcd3b0bcf7/Analysis/GeneralizedRiemannHypothesis/Lean/grh/src/#1}{#2}}
\newcommand{\propdir}{\href{https://github.com/VladimirReshetnikov/ProveIt/tree/9050b04000647d588d0c78d945180abcd3b0bcf7/Analysis/GeneralizedRiemannHypothesis/docs/proposals}{proposal directory}}
\title{\color{ink}\bfseries Direct module duality and Cayley--Hamilton\\in the first-order Ribet argument\\[5pt]\large\normalfont A replacement for two algebraic proof layers in the Lean proof of Dirichlet GRH}
\author{Research note for ProveIt}\date{27 November 2026}
\begin{document}\maketitle
\begin{abstract}
The first-order congruence argument of the readable Alp\"oge--Claude Lean development contains a three-stage Fitting-ideal estimate, expressed using modules of functions and successive enlargements of a relation module. We identify the function module with the exact dual of the finite torsion module generated by its evaluation vectors. This permits an elementary proof of the same triangular Fitting-ideal theorem: choose elementary-divisor generators once, add the prescribed relations, and lift a single determinant from the specialization at zero. The cost of the final step is exactly the length of the $C$-torsion in that specialization, rather than its coarser bound by the number of generators times $v(C)$. A second, independent replacement applies Cayley--Hamilton directly to a quotient on which a formal variable acts by right multiplication, dispensing with the intermediate polynomial adjugate and its evaluation lemma. Both preserve the existing public interfaces. This is a mathematical refactoring proposal, not a new independent proof of GRH or a compiled Lean patch.
\end{abstract}
\section{Recommendation and scope}
The recommended target is \src{New/Congruences/FirstOrder/FittingBlock.lean\#L137}{\code{exists_mul_prod_X_mem_fittingIdeal}}, called in \src{New/Congruences/FirstOrder/RibetMatrices.lean\#L769}{RibetMatrices, line 769}. The other target is \src{New/Congruences/FirstOrder/RibetMethod.lean\#L216}{\code{exists_ribet_polynomial}}, called at line 534 of the same file. These are inside the construction of the common Ribet element that clears theta-lift projection denominators. The arithmetic inputs---the tied Selmer bounds, local crystalline defects, labels, and odd-class bounds---are unchanged. The determinant lower bound, main inequality and zero-exclusion argument are also unchanged.

The distinction from the planned flag-row and torsion-grid changes matters: those select or combine operators after the congruence estimate, or select the character at which a polynomial is evaluated. Here the polynomial certificate itself is produced by a different proof, before either step. The replacement is compatible with the planned proposals without counting their savings again.

Source snapshot: \href{https://github.com/VladimirReshetnikov/ProveIt/tree/9050b04000647d588d0c78d945180abcd3b0bcf7/Analysis/GeneralizedRiemannHypothesis}{9050b04000647d588d0c78d945180abcd3b0bcf7}. Selected blobs were fetched at this revision. The repository describes the baseline as independently kernel checked; see the \href{https://math.mit.edu/~drew/shortgrh/}{verification record}. This note does not rerun those checks or assume that the new lemmas have been formally verified.
\section{The exact dual of a module of evaluations}
Fix a discrete valuation ring $\OO$ with uniformizer $\pi$, and normalize $v(\pi)=1$. Let $p\ne0$, $A=\OO/(p)$, $m=v(p)$. Write $D(V)=\Hom_{\OO}(V,A)$ for a finite-length $\OO$-module killed by $p$. No finiteness of the residue field is necessary. In particular, ``finite'' below means finite length, not necessarily finite cardinality.

\begin{lemma}[Finite torsion duality]\label{lem:dual}
For $V$ killed by $p$ and finitely generated, $D(V)$ has the same length as $V$, separates points of $V$, and restriction $D(V)\to D(W)$ is surjective for every submodule $W\subseteq V$. Its kernel consists precisely of functionals vanishing on $W$. Thus
\[
 \len(D(V)/\operatorname{Ann}(W))=\len W.
\]
\end{lemma}
\begin{proof}
Write $V\cong\bigoplus_i\OO/(\pi^{a_i})$, where $0\le a_i\le m$. The dual of each summand is $\pi^{m-a_i}A$, isomorphic to the original summand. This proves the length statement and point separation. For extension from $W$, adjoin one generator $u$ at a time. The annihilator of $u+W$ is $(\pi^k)$, $k\le m$. To extend $f$ to $W+\OO u$, choose $t\in A$ with $\pi^k t=f(\pi^k u)$. Such $t$ exists because the right side is killed by $\pi^{m-k}$ and hence belongs to $\pi^k A$. Then $f(w)+bt$ is well-defined on $w+bu$. Finitely many extensions suffice. The restriction map is therefore onto, giving the quotient-length formula. The case $m=0$ has only zero modules.
\end{proof}

Let $J$ be a finite linearly ordered set of variables and $R=\OO[X_j:j\in J]$. Let $s$ be the finite index set of rows, $d=|s|$, $E=R^s$, $N\subseteq E$ a relation module, $\Phi:\OO^n\to E$ an $\OO$-linear map, and $z:G\to\OO^n$ an arbitrary function. Assume $\Phi(p\OO^n)\subseteq N$. For any $Q\supseteq N$, put
\[
 M_Q=E/Q,\quad y_Q(g)=\Phi(z(g))+Q,\quad
 V_Q=\Span_{\OO}\{y_Q(g):g\in G\}\subseteq M_Q.
\]
Here $V_Q$ is a submodule of the finite $p$-torsion image of $\Phi$. We do \emph{not} assume $z(G)$ spans $\OO^n$, and $M_Q$ need not have finite length over $\OO$.

The source \src{Lib/General/Algebra/FittingIdealsAndLengths.lean\#L591}{\code{formsMod}} is the module $Z_Q$ of functions $g\mapsto\Psi(z(g))$ for $\Psi:\OO^n\to A$ vanishing on $\Phi^{-1}(Q)$ (the extra singleton coordinate in Lean is immaterial).
\begin{lemma}[Evaluation-dual identification]\label{lem:eval}
Evaluation induces a canonical $\OO$-linear isomorphism
\[
 D(V_Q)\xrightarrow{\ \sim\ } Z_Q,\qquad f\longmapsto(g\mapsto f(y_Q(g))).
\]
If $B\subseteq G$ and $W_B=\Span_{\OO}\{y_Q(g):g\in B\}$, the functions vanishing on $B$ correspond exactly to $\operatorname{Ann}(W_B)$. If all $c\in Z_Q$ satisfy $c(g)-c(h)=0$, then $y_Q(g)=y_Q(h)$; the assertion also holds for arbitrary finite linear combinations and scalar multiples.
\end{lemma}
\begin{proof}
A source functional factors through the image $U_Q$ of $\Phi$ in $M_Q$. Restriction to $V_Q$ gives the function. Conversely extend a functional on $V_Q$ to $U_Q$ by Lemma~\ref{lem:dual}, then precompose with $\Phi$. Injectivity holds since the evaluation vectors span $V_Q$. The remaining assertions are restriction and point separation.
\end{proof}

This identification is the central simplification: work in the finite module that actually accounts for the functions, rather than transport length comparisons through several larger spaces of linear forms. It remains valid for an infinite set $G$ and does not choose a generating set of $G$.
\section{Transporting determinants without extra relation modules}
For $Q\subseteq R^s$, denote by $\FF(E/Q)$ the ideal generated by determinants of $d$-by-$d$ matrices with rows in $Q$ (the convention of the source). When $d=0$, this is the unit ideal.
\begin{lemma}[One added vector]\label{lem:transport}
If $a\in R$ and $av\in Q$ for $v\in E$, then
\[
 a\FF(E/(Q+Rv))\subseteq \FF(E/Q).
\]
\end{lemma}
\begin{proof}
Write each row of a determinant for $Q+Rv$ as $q_i+b_iv$. Multilinearity leaves the determinant of the $q_i$ and terms with exactly one $v$ row; terms with two copies of $v$ vanish by alternation. Multiplying by $a$ puts every surviving row in $Q$ (for the term with no $v$, multiply any row; equivalently use ideal closure). This proves the inclusion for generators and hence for the ideal. The case $d=0$ is immediate.
\end{proof}
This is the elementary transport already implicit in \src{Lib/General/Algebra/FittingIdealsAndLengths.lean\#L494}{\code{mul_det_mem_fittingIdeal}}. We keep it. The new part is deriving the entire cost of the first enlargement directly from the elementary divisors of a single finite torsion module, and deriving the final cost by specialization. This avoids the general induction on forms that currently supplies both steps.
\section{The triangular Fitting-ideal replacement}
\subsection{Exact hypotheses}
We restate the source theorem. Let $I\subseteq J$, let $C\in\OO$, and let $\tau_j,\tau'_j,\phi_j,\phi'_j\in G$. In addition to $p\ne0$ and the preceding setup, assume:
\begin{enumerate}
\item For each $g\in s\subseteq G$, $y_N(g)$ equals the image of the corresponding standard basis vector $e_g$. Thus $pM_N=0$ as an $R$-module, even though it may have infinite $\OO$-length.
\item In $M_N$, $y_N(\tau_j)=X_jy_N(\phi_j)$ and $y_N(\tau'_j)=X_jy_N(\phi'_j)$ for every $j$.
\item If $j\in I$ and $i<j$ with $i\notin I$, then $X_iM_N=0$.
\item There exists $B\subseteq G$ and a cost bound $b\in\mathbb N\cup\{\infty\}$ with the following two properties. Write $Z=Z_N$, $Z_B=\{c\in Z:c|_B=0\}$. For $j\in I$ and $c\in Z_B$, if $c(\tau_i)=c(\tau'_i)$ for all $i<j$ and $c(\phi_i)=c(\phi'_i)$ for $i<j$ in $I$, then
\[
 C(c(\tau_j)-c(\tau'_j))=0.
\]
For every submodule $S\subseteq Z_B$ vanishing at all $\tau_j,\tau'_j$ and satisfying $c(\phi_j)=c(\phi'_j)$ for $j\in I$,
\[
 \len(Z/Z_B)+\len(S/S[C])\le b,
 \qquad S[C]=\{c\in S:Cc=0\}.
\]
\end{enumerate}
These are exactly the basis, tie, initial-variable and \src{New/Congruences/FirstOrder/FittingBlock.lean\#L46}{\code{HasTriangleCost}} conditions. We do not strengthen the Selmer hypotheses.
\begin{theorem}[Direct certificate]\label{thm:main}
Under these assumptions with $C\ne0$, there is a polynomial $F\in R$ such that
\[
 F\prod_{j\in I}X_j\in\FF(E/N),\qquad
 v(F(0))\le b+|I|v(C)+\len(S_*[C])
 \le b+(|I|+d)v(C),
\]
where $S_*$ is a submodule satisfying the cost hypothesis, defined explicitly below. In particular, the original Lean output follows without changing its statement.
\end{theorem}
\begin{proof}
\textbf{Step 1: absorb all detection points at their exact length cost.}
Let $W_B\subseteq V_N$ be the span of $y_N(g)$ for $g\in B$. By Lemma~\ref{lem:eval} and duality,
\[
 e_1=\len W_B=\len(Z/Z_B).
\]
Choose elementary-divisor generators $w_1,\dots,w_h$ of $W_B$, of orders $\pi^{a_1},\dots,\pi^{a_h}$, and lifts $\tilde w_i\in E$. Add them to $N$ over $R$, defining $K=N+\sum_iR\tilde w_i$. Since $\pi^{a_i}\tilde w_i\in N$, repeated application of Lemma~\ref{lem:transport} gives
\[
 \pi^{e_1}\FF(E/K)\subseteq \FF(E/N),\qquad e_1=\sum_i a_i.
\]
Every function in $Z_Q$ for $Q\supseteq K$ vanishes on $B$: all $y_N(g)$ for $g\in B$ have become zero. No assumption that $B$ is finite is required.

\textbf{Step 2: use the triangular relations in their given order.}
For each $j\in I$, set $v_j=\Phi(z(\phi_j)-z(\phi'_j))$. Let $Q_j$ be $K$ enlarged by $Rv_i$ for $i\in I$, $i<j$. For every $i<j$, $X_i v_i=0$ in $M_{Q_j}$: for $i\in I$ this follows from the enlargement, and otherwise from hypothesis 3. By the ties, all functions in $Z_{Q_j}$ satisfy the previous $\tau$ equalities and the requisite previous $\phi$ equalities. They vanish on $B$, so the triangle condition gives
\[
 C(c(\tau_j)-c(\tau'_j))=0\quad(c\in Z_{Q_j}).
\]
By point separation in Lemma~\ref{lem:eval}, $C(y_{Q_j}(\tau_j)-y_{Q_j}(\tau'_j))=0$. The ties now show $CX_jv_j\in Q_j$. Apply Lemma~\ref{lem:transport} while adding $v_j$. After the last step put $T=K+\sum_{j\in I}Rv_j$. We have
\[
 C^{|I|}\left(\prod_{j\in I} X_j\right)\FF(E/T)\subseteq\FF(E/K).
\]

\textbf{Step 3: specialize once, and lift a determinant.}
Let $\mathfrak x=(X_j:j\in J)$, and put $T_0=T+\mathfrak x E$. Then
\[
 M_0=E/T_0\simeq \OO^s/\overline T,
 \qquad \overline T=\{t(0):t\in T\}.
\]
By the basis condition $p\OO^s\subseteq\overline T$, and the images of $y_{T_0}(g)$ for $g\in s$ span all of $M_0$ over $\OO$. Thus $V_{T_0}=M_0$, and $S_*:=Z_{T_0}\simeq D(M_0)$. This $S_*$ is contained in $Z_B$; every $\tau$ and $\tau'$ value vanishes (all variables are zero); the $\phi$ equalities for $I$ hold by construction of $T$. The cost assumption applies.

Choose an $\OO$-basis of the full-rank lattice $\overline T$ in $\OO^s$. Its determinant has valuation $\len M_0$ by elementary divisors. Lift its $d$ rows to $T$, giving a square matrix $H(X)$, and set $D(X)=\det H(X)$. Then $D\in\FF(E/T)$ and
\[
 v(D(0))=\len M_0=\len S_*=\len(S_*/S_*[C])+\len(S_*[C]).
\]
Since $M_0$ is generated by $d$ elements, it has at most $d$ cyclic elementary divisors, say of lengths $a'_i$. Hence
\[
 \len(S_*[C])=\sum_i\min(a'_i,v(C))\le d\,v(C).
\]
Finally set $F=\pi^{e_1}C^{|I|}D$. The two transport inclusions give the required ideal membership. Its constant coefficient satisfies
\begin{align*}
 v(F(0))&=\len(Z/Z_B)+|I|v(C)+\len S_*\\
 &\le b+|I|v(C)+\len(S_*[C]),
\end{align*}
which proves both claimed bounds.
\end{proof}
\paragraph{Degenerate cases.} If $C=0$ and $I\ne\varnothing$, the original (extended-natural-valued) bound is automatic with $F=0$. If $I=\varnothing$, no triangle multiplier is needed: the argument gives $F=\pi^{e_1}D$, $v(F(0))=e_1+\len S_*$ even for $C=0$. If $d>0$, the old coarse cost $d v(0)$ is infinite; if $d=0$, $E=0$ and $F=1$ works. If $b=\infty$, the displayed finite construction still works for $C\ne0$, although the old bound is vacuous. These conventions should be handled explicitly in the Lean wrapper.

\paragraph{Quantitative gain.} The exact saving against the old last-stage allowance is
\[
 d v(C)-\sum_i\min(a'_i,v(C))\ge0.
\]
If the specialization needs only $h'<d$ generators this is at least $(d-h')v(C)$. A universal sharpening of the downstream global exponent is not claimed, since the invariant factors of $M_0$ depend on the actual Galois relations. The main gain is the simpler proof, not a numerical change in the GRH statement.
\subsection{An explicit certificate with a strict saving}
Work over $\OO=\mathbb Z_{(3)}$ with $p=3^4$, $C=3^2$, $R=\OO[X]$, $E=R^3$, and $I=\{1\}$. Take
\[
 N=\langle p e_1,p e_2,p e_3,e_3-Xe_2,Ce_3\rangle_R.
\]
Choose evaluation vectors $e_1,e_2,e_3,3e_1,0$, with the detection set evaluating to $3e_1$ and the tie evaluating to $y(\tau)=e_3$, $y(\phi)=e_2$, $y(\tau')=y(\phi')=0$. One can take $\Phi$ to be the inclusion of constant vectors. The evaluation module is $\OO/(p)\oplus\OO/(p)\oplus\OO/(C)$: after eliminating $e_3$ the only additional relation is $CXe_2=0$. The detection submodule has length $3$. The triangle holds because $Ce_3=0$ in the quotient. The final permitted functions have their second and third values zero, and their first value killed by $3$. Thus the cost hypothesis holds with $b=3$.

Add $3e_1$, then $e_2$. The specialization has just one cyclic factor of length $1$. Lift its row basis by the matrix
\[
 H(X)=\begin{pmatrix}3&0&0\\0&1&0\\0&-X&1\end{pmatrix},\qquad F=3^3 C\det H=3^6.
\]
Indeed $FX$ is the determinant of the matrix with rows $p e_1$, $Ce_3$, $e_3-Xe_2$, all in $N$. The old bound is $3+(1+3)\cdot2=11$, while the new bound is $3+2+1=6$, attained by this certificate. This illustrates the exact torsion-length improvement; it does not predict a five-unit saving for the actual arithmetic application.

\section{Why the specialization cannot be moved earlier}
The order of operations is essential. For example, in $R=\OO[X]$, let $E=R$, $N=(p,X)$, $p=\pi^m$. Then $E/N\cong\OO/(p)$, and the specialized dual has length $m$. But the Fitting ideal in $R$ is $(p,X)$, not merely $(p)$. Conversely, consider $N=(p)$: the specialized module is again $\OO/(p)$, while $E/N=(\OO/p)[X]$ has infinite $\OO$-length. One must take the dual of $V_Q$, not of $M_Q$, during the first two steps. The final specialization is only used to choose lifts in $T$; replacing $T$ by $T_0$ in the Fitting transport would lose precisely the polynomial information that the Ribet coefficient requires.

Another failure mode is defining the dual of all of $\OO^n/\Phi^{-1}Q$ and silently identifying its length with the length of $Z_Q$. For instance, if $z(G)$ uses only the first basis vector of $\OO^2$ and both directions have order $p$, then the full dual has length $2m$, but $Z_Q$ has length $m$. Restriction to $V_Q$ is the necessary correction. The extension lemma proves that no information needed for the functions is lost.
\section{Cayley--Hamilton directly on the Ribet quotient}
The second proposal removes a different intermediate construction. It does not require Theorem~\ref{thm:main}. We follow the input of \src{New/Congruences/FirstOrder/RibetMethod.lean\#L216}{\code{exists_ribet_polynomial}}. Let $R_0,L$ be commutative rings, $\Lambda$ an $R_0$-algebra (possibly noncommutative), and $\rho:\Lambda\to M_d(L)$ an $R_0$-algebra map. Suppose all characteristic-polynomial coefficients of $\rho(w)$ are in the image of $R_0$; the source derives this from its spanning hypothesis using Amitsur. Let $\mathfrak a\subseteq\Lambda$ be a two-sided ideal with $\ker\rho\subseteq\mathfrak a$. Let $y=(y_i)_{i\in K}$ be a finite column, and $I_y=\sum_i\Lambda y_i$. Suppose a matrix $B$ has rows satisfying
\[
 (By)_i\in \mathfrak a+I_y\rho_0
\]
for a fixed $\rho_0\in\Lambda$. In the quotient $Q=\Lambda/\mathfrak a$, write this as $By=B_1y\rho_0$ for some $B_1\in M_K(\Lambda)$.
\begin{theorem}[Direct Ribet polynomial]\label{thm:ch}
There is $P\in R_0[X]$ whose constant term maps to $\det(\rho(B))$ (with blocks flattened) and for which $zP(\rho_0)\in\mathfrak a$ for all $z\in I_y$. The two trace conclusions of the source follow unchanged.
\end{theorem}
\begin{proof}
The central variable in $\Lambda[X]$ acts on $Q$ by right multiplication by $\overline{\rho_0}$. Each coefficient $w\in\Lambda$ acts by left multiplication by $\bar w$. These actions commute, giving a ring homomorphism $\Lambda[X]\to\operatorname{End}_{R_0}(Q)$. On $Q^K$ use its matrix extension. Put $D_X=B-XB_1$. By the assumed relation, $D_X\bar y=0$ for this action.

Extend $\rho$ coefficientwise to $\rho_X$ and flatten blocks, obtaining a representation $\varphi_X:M_K(\Lambda[X])\to M_D(L[X])$, $D=d|K|$. Amitsur's formula implies that all coefficients of the characteristic polynomial of $\varphi_X(D_X)$ lie in the image of $R_0[X]$ (as in the source). Choose lifts $c_i(X)\in R_0[X]$ and write its lifted monic characteristic polynomial as
\[
 H(t)=t^D+c_{D-1}(X)t^{D-1}+\cdots+c_0(X).
\]
By Cayley--Hamilton, $\varphi_X(H(D_X))=0$. Every coefficient of every matrix entry of $H(D_X)$ therefore belongs to $\ker\rho\subseteq\mathfrak a$. It follows that $H(D_X)$ acts by zero on $Q^K$. Since all positive powers of $D_X$ annihilate $\bar y$, this gives $c_0(X)\bar y=0$. In terms of the action, $\bar y_i c_0(\overline{\rho_0})=0$ for every $i$. Take $P(X)=(-1)^D c_0(X)$; its constant term maps to $\det(\rho(B))$ by specialization at $X=0$. Left linear combinations give the assertion for $I_y$. If $K$ is empty, take $P=1$; the assertion is then vacuous.

For the trace assertions, $\rho(\rho_1 P(\rho_0))$ is the image of an element of $\Lambda$, so its trace is in the image of $R_0$ by the characteristic coefficient hypothesis (the source assumes $d>0$). If $z_0\in I_y$ and $\rho(z_0)=\delta I_d$, then $a=\rho_1 z_0 P(\rho_0)\in\mathfrak a$, and its trace under $\rho$ is $\delta\,\operatorname{tr}(\rho(\rho_1P(\rho_0)))$. These are exactly the two remaining conclusions.
\end{proof}
The point is that we never need an identity of the form $B'B_X=P I$ in a polynomial matrix algebra. The only consumer wants a scalar annihilating $\bar y$. The characteristic polynomial already provides this, without choosing $B'$ or proving a separate evaluation identity. The use of right multiplication, rather than substituting $\rho_0$ into a noncommutative polynomial ring as a ring map, is essential.

\subsection{A smaller Amitsur lifting lemma}
The source \src{New/Congruences/FirstOrder/RibetMethod.lean\#L51}{block Amitsur lemma} uses matrices with one nonzero block row and a rectangular product identity. A simpler spanning family is the set of matrices with a single nonzero block $w$ at $(i,j)$. If $i=j$, the flattened characteristic polynomial is $t^{D-d}\operatorname{charpoly}(\rho(w))$. Otherwise reorder the block indices so $i<j$; the matrix is strictly block upper triangular and its characteristic polynomial is $t^D$. These statements hold over an arbitrary commutative coefficient ring, even nonreduced rings, because they use triangularity rather than an inference from nilpotence. The elementary blocks span the whole matrix algebra, so the existing Amitsur span lemma proves the same interface. No change to Amitsur's general theorem is required. The proof over $R_0[X]$ uses the same coefficientwise representation and monomial span as the original.
\section{Integration and expected effect}
The readable source has about 300,000 code lines according to its \href{https://github.com/VladimirReshetnikov/ProveIt/blob/9050b04000647d588d0c78d945180abcd3b0bcf7/Analysis/GeneralizedRiemannHypothesis/Lean/grh/README.md}{README}. The proposals in this note target a small but conceptually nontrivial part, not an order-of-magnitude global reduction. The following are \emph{regions to replace or audit}, not measured net savings.
\begin{center}\small
\begin{tabularx}{\linewidth}{@{}p{.29\linewidth}X@{}}\toprule
Source region & Proposed action \\\midrule
\src{New/Congruences/FirstOrder/FittingBlock.lean\#L65-L209}{FittingBlock, 65--209} & Replace the constant-coefficient helper and triangular proof with evaluation-dual identification, elementary-divisor transport and specialization. Preserve the public triangular theorem.\\
\src{Lib/General/Algebra/FittingIdealsAndLengths.lean\#L591-L721}{FittingIdealsAndLengths, 591--721} & Keep formsMod and antitonicity. Audit the generic strict-descent/relative-length helpers (especially lines 635--721) for deletion after the replacement. Keep shared Fitting definitions and single-vector transport.\\
\src{New/Congruences/FirstOrder/RibetMethod.lean\#L51-L280}{RibetMethod, 51--280} & Simplify block Amitsur and replace the adjugate route. Keep the exported Ribet polynomial theorem and its trace conclusions.\\
\src{New/Congruences/FirstOrder/RibetMatrices.lean\#L769}{RibetMatrices, 769} & Existing call to the triangular theorem should be unchanged.\\
\src{New/Congruences/FirstOrder/RibetElement.lean\#L289}{RibetElement, 289, 404} & Retain the common element and coefficient valuation.\\
\src{New/Congruences/FirstOrder/DenominatorBound.lean\#L1321}{DenominatorBound, 1321} & Retain the exact denominator-clearing interface.\\\bottomrule
\end{tabularx}
\end{center}
The three inspected regions span about 500 physical lines including comments, declarations retained, and shared helper code. The actual Lean savings may be substantially smaller, especially if the torsion dual extension theorem requires additional formal infrastructure. The principal qualitative improvement is to collapse two function-module enlargements into one duality identification and one determinant specialization, and to remove the adjugate certificate entirely. It would be misleading to count all lines in these regions as deleted or to extrapolate to the entire project.

The full integration is intentionally separated into two independently testable patches. A developer should first port Theorem~\ref{thm:main} and rebuild the callers, then separately port Theorem~\ref{thm:ch}. See \path{LEAN_PLAN.md} for exact proof obligations and acceptance gates. The source audit contains immutable blob hashes and the local verification script. No repository mutation was made.
\section{Nonduplication against the current archive}
The \propdir{} contains 18 subdirectories at the pinned commit. Their main mechanisms are grouped below. The closest predecessor is the flag-row construction: it combines clearing operators on the lift space, not the evaluation-module certificate from which the individual congruence bound is obtained. The other potentially similar method, discrete Hecke's use of Smith form, estimates exceptional local Hecke directions, not Selmer functions or Fitting coefficients.
\begin{center}\small
\begin{tabularx}{\linewidth}{@{}p{.48\linewidth}X@{}}\toprule
Proposal families & Distinct target \\\midrule
Centered moments; cyclotomic interpolation; sparse theta; shorter theta; proof simplifications; RH simplifications & Spanning, cutoff, zero detection or choice of torsion twist; after or outside the certificate construction.\\
CM norm; signed theta gains & Local sharpness and divisibility accounting.\\
Discrete Hecke; edge Hecke & Local elementary divisors and integral orbit traces.\\
Gram adjoint; orbital theta exchange & Comparison operator and exchange identities.\\
Finite-quotient class count & Counting unitary classes, not Selmer relation modules.\\
Order tests; two order tests; robust labels & Local good-type projector arithmetic.\\
Newton torsion & Choosing a common primitive torsion order.\\
Theta-lift simplifications & Flag-row projection clearing, labels, explicit sparse places.\\\bottomrule
\end{tabularx}
\end{center}
The audit is at directory, abstract/README, and targeted report-text level. We retrieved all 18 tree entries, read their main summaries and selected relevant full sections; all 28 listed TeX sources were screened for the relevant theorem names and keywords in two successful batches (the initial larger request hit the connector call limit). This is not an exhaustive search of all possible paraphrases in every file, nor a claim of literature priority. See \path{proposal_audit.tsv} for per-folder scope. The mathematical ingredients (DVR duality, Fitting transport and Cayley--Hamilton) are standard; the proposal is their direct assembly at these precise interfaces.
\section{Verification and next research questions}
The included exact-arithmetic Python tests check finite torsion duality (including failure of the full-dual shortcut), determinant transport, a concrete triangular certificate with a strict saving, and direct Cayley--Hamilton on noncommutative matrix-coefficient examples. These tests are regressions, not evidence that all abstract statements are proved by enumeration. The universal arguments are above.

No Lean executable or full project cache was available in the working environment. The original upstream proof was not rebuilt. Neither replacement has been integrated or kernel checked. The report was compiled from the included TeX; the recorded checks and source hashes are provided for reproduction.

Further research: (1) Can $\len S_*[C]$ be controlled by the dimension of the actual tie-evaluation span rather than by $d$ in the arithmetic caller? (2) Can the finite evaluation-dual identification be reused in the crystalline-colength proof without changing any local hypotheses? (3) Can the common Ribet element be constructed directly over a multilinear coefficient quotient while preserving ideal membership and the exact coefficient of $\prod X_j$? The last question is not answered here; merely quotienting by $X_j^2$ can destroy polynomial ideal information and must be checked at all callers.
\end{document}
